\hypertarget{chapter-5-build-systems-and-continuous-compilation}{%
\chapter{Build Systems and Continuous Compilation}\label{chapter-5-build-systems-and-continuous-compilation}}

The previous two chapters established how a large project's source should be organized and tracked. This chapter is about the step that turns that source into a finished PDF: running the compiler enough times, in the right order, with the right auxiliary tools, to produce correct output, and doing so with as little manual intervention as the project's size demands. A single-page document can be built by typing \texttt{lualatex\ file.tex} and glancing at the result. A book with a bibliography, an index, and cross-references scattered across forty chapter files cannot, not because any individual step is hard, but because the sequence of steps has to be run correctly and in order every time, and a human doing that by hand will eventually skip a step, run it out of order, or forget it entirely under deadline pressure.

\hypertarget{the-multi-pass-problem}{%
\section{The Multi-Pass Problem}\label{the-multi-pass-problem}}

LaTeX's compiler processes a document from beginning to end in a single pass, but several features a book-length document depends on cannot be resolved in a single pass in principle, not merely as an implementation limitation. A table of contents needs to know every section's title and page number before it can be typeset, but the table of contents itself typically appears near the beginning of the document, before most of those sections have been processed. A cross-reference to ``see Chapter 12'' needs to know what page Chapter 12 ends up on, which is not settled until the entire document, including everything after the reference, has been laid out.

LaTeX resolves this by writing what it learns during one pass into an auxiliary file, the \texttt{.aux} file, and reading that file back in at the start of the next pass. On the first pass, every \texttt{\textbackslash{}label} records its current page and section number into the \texttt{.aux} file, but every \texttt{\textbackslash{}ref} to a label defined later in the document has nothing to read yet, since that label has not been processed and its position has not been recorded, so it renders as a placeholder question mark. On the second pass, the \texttt{.aux} file now contains every label's position from the first pass, so every reference resolves correctly, assuming nothing in the document has shifted enough between passes to invalidate what the first pass recorded. In practice, a document with many cross-references sometimes needs a third pass, since a reference resolving on the second pass can itself shift page breaks enough to move a label from the first pass, and only stabilizes once a full pass makes no further changes.

Bibliographies and glossaries follow the same pattern with an extra step: BibTeX or Biber reads the citations recorded by a first LaTeX pass, produces a formatted bibliography, and that bibliography then needs to be pulled back into the document on a subsequent LaTeX pass, meaning a document with a bibliography typically needs LaTeX, then the bibliography tool, then LaTeX again at minimum, and often once more to let cross-references to bibliography entries settle. An index built with \texttt{makeindex} and a glossary built with the \texttt{glossaries} package's own processing tool add further instances of the same pattern: something outside LaTeX itself has to run between two LaTeX passes to produce a resource the second pass consumes.

\hypertarget{automating-the-sequence}{%
\section{Automating the Sequence}\label{automating-the-sequence}}

None of this sequence is complicated to describe, but a sequence of five or six commands, run in the correct order, with the correct arguments, every time any file in the project changes, is exactly the kind of task that should not be left to memory. \Pgm{latexmk} exists to solve this directly: given a top-level \texttt{.tex} file, it inspects what auxiliary tools the document actually needs (by watching for \texttt{.bib}, \texttt{.idx}, or glossary files it needs to process), runs the correct sequence of compiler and auxiliary-tool invocations automatically, and, critically, keeps rerunning until the output stabilizes rather than running a fixed number of passes regardless of whether they were sufficient. A project set up with a \texttt{latexmkrc} file specifying the engine (LuaLaTeX, per this book's conventions) and any non-default options needs nothing more than \Pgm{latexmk} typed at the top level to produce a correct, fully resolved PDF, regardless of how many auxiliary tools the current draft happens to require.

\Pgm{arara} takes a different approach to the same problem, driven by directives written as comments inside the \texttt{.tex} file itself, specifying which tools to run for that particular document. This is useful when different documents in the same project need different build sequences, or when the build sequence needs to be readable by someone opening the source file itself rather than a separate configuration file. For a single coherent book project with one consistent build sequence throughout, \Pgm{latexmk}'s external configuration is usually the better fit; for a collection of loosely related documents with varying needs, \Pgm{arara}'s per-file directives can be the more maintainable choice.

A Makefile is worth adding on top of either tool once a project has multiple build targets: a full book, an excerpt for review, a version with tracked-changes markup visible, a print-ready PDF with different margins than the screen version. Encoding these as named Make targets, each invoking \Pgm{latexmk} or \Pgm{arara} with the appropriate options, turns ``how do I build the review copy again'' into \texttt{make\ review}, which is a small convenience for a single author working alone and a considerable one for a project anyone else needs to build without first reconstructing the author's usual sequence of commands from memory.

\hypertarget{continuous-compilation-during-drafting}{%
\section{Continuous Compilation During Drafting}\label{continuous-compilation-during-drafting}}

The tooling described so far handles producing a final, fully resolved PDF, but a full multi-pass build with bibliography and glossary processing is more than is needed, and often too slow to be pleasant, while actively drafting a single chapter. \Pgm{latexmk}'s \texttt{-pvc} (preview continuously) mode addresses this directly: it watches the project's files for changes and triggers a rebuild automatically on save, paired with a PDF viewer that reloads the output, so that the feedback loop between writing a sentence and seeing it typeset shrinks to the time a single build takes rather than requiring a manual command after every edit.

This matters more for a project organized the way Chapter 3 describes than it would for a single-file document, because a chapter-scoped edit only needs to rebuild quickly if the build tooling is actually taking advantage of that scoping. Combining \texttt{\textbackslash{}includeonly} with a continuous build restricted to the chapter currently being drafted keeps the feedback loop fast even in a project whose full build, bibliography and all, takes long enough to be disruptive if triggered on every keystroke-adjacent save. The full build remains necessary before anything is considered finished, since a chapter-scoped build will not catch a broken cross-reference into a chapter that was excluded, but it does not need to run on every edit, only at the points where the document actually needs to be checked as a whole.

With the project organized, tracked, and buildable on demand, the next part of the book turns to the layer above all of this infrastructure: the macros and interfaces an author actually writes with, and how to design them so that a project's growing complexity is managed by the tooling rather than by the author's memory.
